% ====================================================================
% VİZE RAPOR ŞABLONU – Türk Masalları Derlemesi Web Sayfası
% --------------------------------------------------------------------
% DERLEME: Overleaf'te Menu > Compiler > XeLaTeX (pdfLaTeX de çalışır).
%
% KURALLAR (şablonu kullanmadan önce okuyun):
%  1. Bölüm başlıklarını SİLMEYİN ve sırasını DEĞİŞTİRMEYİN.
%  2. [KÖŞELİ PARANTEZ] içindeki ve gri renkli açıklama metinlerini
%     kendi bilgilerinizle değiştirin; açıklama metinlerini silin.
%  3. MEBİS'e YALNIZCA bu raporun PDF'ini yükleyin.
%     ZIP ve HTML dosyası YÜKLEMEYİN. Tüm kontroller GitHub'dan yapılır.
%  4. Kapaktaki GitHub ve YouTube bağlantıları doğru ve tıklanabilir
%     olmalıdır. Teslimden önce PDF'teki bağlantılara tıklayıp deneyin.
%  5. Haftalık katkı tablosundaki görseller için resim dosyalarını
%     (hafta1.png, hafta2.png, ...) Overleaf projesine yükleyin.
%     Dosya adlarında boşluk ve Türkçe karakter kullanmayın.
%     Resim yoksa tabloda "görsel ekleyin" kutusu görünür.
%  6. Yapay zekâ kullandıysanız Bölüm 5'i doldurun.
% ====================================================================
\documentclass[12pt,a4paper]{article}

\usepackage{iftex}
\ifPDFTeX
  \usepackage[T1]{fontenc}
  \usepackage[utf8]{inputenc}
\else
  \usepackage{fontspec}
\fi
\IfFileExists{turkish.ldf}{\usepackage[turkish]{babel}}{\usepackage[provide=*,turkish]{babel}}
\usepackage[margin=2.4cm]{geometry}
\usepackage{setspace}
\usepackage{graphicx}
\usepackage{booktabs}
\usepackage{longtable}
\usepackage{array}
\newcolumntype{P}[1]{>{\raggedright\arraybackslash}p{#1}}
\usepackage{colortbl}
\usepackage{xcolor}
\usepackage{enumitem}
\usepackage{listings}
\usepackage{amssymb}
\usepackage{caption}
\usepackage{fancyhdr}
\usepackage{hyperref}

\setlength{\headheight}{15pt}
\setlength{\emergencystretch}{3em}
\onehalfspacing
\setlength{\parindent}{0pt}
\setlength{\parskip}{6pt}
\setlist{itemsep=2pt, topsep=4pt}
\hypersetup{colorlinks=true, linkcolor=blue!50!black, urlcolor=blue!50!black}

\definecolor{ana}{HTML}{14325C}
\definecolor{acik}{HTML}{EEF2F8}

\lstset{basicstyle=\ttfamily\small, backgroundcolor=\color{acik},
  frame=single, rulecolor=\color{ana}, breaklines=true,
  columns=fullflexible, keepspaces=true, showstringspaces=false}

\pagestyle{fancy}
\fancyhf{}
\fancyhead[L]{Masalsı Pırıltı}
\fancyhead[R]{Vize Raporu}
\fancyfoot[C]{\thepage}

% ---------- Yardımcı komutlar ----------
% Gri açıklama metni (silinecek yönergeler için)
\newcommand{\aciklama}[1]{{\color{gray}\itshape (#1)}}

% Haftalık tablo görseli: dosya yoksa "görsel ekleyin" kutusu çıkar
\newcommand{\haftagorsel}[1]{%
  \IfFileExists{#1}{%
    \includegraphics[width=\linewidth,height=4.2cm,keepaspectratio]{#1}%
  }{%
    \fbox{\parbox[c][2.6cm][c]{\dimexpr\linewidth-2\fboxsep-2\fboxrule\relax}{%
      \centering\small\color{gray} Ekran görüntüsü / görsel ekleyin\\ \texttt{#1}}}%
  }%
}
% Haftalık katılım kutuları (işaretlemek için $\square$ yerine $\boxtimes$ yazın)
\newcommand{\katilimkutulari}{\par\vspace{3pt}\small $\square$ Ders içinde katılım sergiledim\par $\boxtimes$ Telafi yaptım}

% Ek görsel: dosya yoksa hiçbir şey göstermez
\newcommand{\haftaek}[1]{%
  \IfFileExists{#1}{\\[4pt]\includegraphics[width=\linewidth,height=4.2cm,keepaspectratio]{#1}}{}%
}

% ---------- Kapak bilgileri ----------
\newcommand{\universite} {İSTANBUL MEDİPOL ÜNİVERSİTESİ / GÖRSEL
İLETİŞİM TASARIMI  }  
\newcommand{\dersadi}{PROGRAMLAMAYA GİRİŞ}
\newcommand{\hocaadi}{MERAL ŞENTÜRK}
\newcommand{\ogrenciadi}{İREM AKPINAR}
\newcommand{\ogrencino}{86240029}
\newcommand{\projeadi}{Masalsı Pırıltı}
% Bağlantıları tam adres olarak yazın (https:// ile başlayarak)
\newcommand{\githublink}{\url{https://github.com/ireem1622-ops/masalsi-pirilti}}
\newcommand{\youtubelink}{\url{https://www.youtube.com/watch?v=VIDEO-KIMLIGI}}

\begin{document}

% ====================================================================
% KAPAK
% ====================================================================
\hypersetup{pageanchor=false}
\begin{titlepage}
\centering
\vspace*{1cm}
{\large \universite \par}
\vspace{1.8cm}
{\Huge\bfseries\color{ana} \projeadi \par}
\vspace{0.4cm}
{\Large Türk Masalları Derlemesi Web Sayfası \par}
{\large Vize Raporu \par}
\vspace{1.8cm}
{\large \textbf{Ders:} \dersadi \par}
{\large \textbf{Öğretim Üyesi:} \hocaadi \par}
\vspace{1.2cm}
{\large \textbf{Hazırlayan:} \ogrenciadi \par}
{\large \textbf{Öğrenci No:} \ogrencino \par}
\vspace{1.8cm}
\fbox{\parbox{0.88\linewidth}{\centering\vspace{6pt}
{\large\bfseries Proje Bağlantıları}\\[8pt]
\textbf{GitHub deposu:}\\ \githublink \\[8pt]
\textbf{YouTube sunum videosu:}\\
\url{https://www.youtube.com/watch?v=1Jkv4_qK31A}
\vspace{6pt}}}
\vfill
{\large 6 Kasım 2026 \par}
\end{titlepage}
\hypersetup{pageanchor=true}

\tableofcontents

% ====================================================================
\section{Giriş}
% ====================================================================

Masalsı Pırıltı, sekiz masalı görsellerle bir araya getiren etkileşimli bir masal okuma web sitesidir. Projeyi, ziyaretçilerin masalları kolayca keşfedebileceği, farklı hikâyeler arasında seçim yapabileceği ve okuma görünümünü kendi tercihine göre ayarlayabileceği bir ortam oluşturmak amacıyla hazırladım. Masalları görsel, başlık ve kısa açıklama içeren kartlarla sunarak okuyucunun metnin tamamını açmadan önce hikâye hakkında fikir edinmesini ve ilgisini çeken masala yönelmesini amaçladım. Mavi, mor ve turkuaz tonlarını orman ve geceyi çağrıştıran illüstrasyonlarla birleştirdim, küçük hareketli pırıltılarla sitenin adını destekleyen bir atmosfer oluşturdum. Sitede sekiz masala yer verdim ve sayfanın yapısını, görsellerini, masal metinlerini ve etkileşimlerini tek HTML dosyasında topladım. İçsel CSS ile renkleri, yazı tiplerini, kartların yerleşimini ve farklı ekran genişliklerine uyum sağlayan düzeni belirledim. İçsel JavaScript ile başlığa göre arama, tema filtreleme, rastgele masal seçimi, masalın okuma ekranında açılması, yazı boyutunun değiştirilmesi ve gece moduna geçiş gibi özellikleri gerçekleştirdim.

\subsection{Proje Özeti}
Masalsı pırıltı,sekiz masalı illüstrasyonlarla bir araya getiren etkileşimli bir masal okuma web sitesidir.Projenin amacı,ziyaretçilere masalsı bir atmosfer içinde kolay ve rahat bir okuma deneyimi sunmaktır.Sitede masal adına göre arama,temaya göre filtreleme,rastgele bir masal seçme, masalı okurken isteğe göre yazı boyutunu değiştirme ve gece moduna geçme gibi özellikler bulunmaktadır.Tasarımda renk paleti olarak lacivert,mor ve turkuaz tonları seçilmiştir.Ayrıca sitenin masalsı atmosferini yansıtmak amacıyla küçük hareketli ışıklar kullanılmıştır.

\subsection{Gereksinimlerin Karşılanması}
\begin{center}
\rowcolors{2}{acik}{white}
\begin{tabular}{@{}p{8.2cm}c@{}}
\toprule
\rowcolor{ana}\textcolor{white}{\textbf{Gereksinim}} & \textcolor{white}{\textbf{Durum}} \\
\midrule
Masal sayısı 8--12 arasında (8 masal)         & $\boxtimes$ \\
Tek HTML dosyası                                    & $\boxtimes$ \\
Stil kodları içsel (\texttt{<style>} içinde)        & $\boxtimes$ \\
Betik kodları içsel (\texttt{<script>} içinde)      & $\boxtimes$ \\
İnternetsiz açılıp çalışıyor                        & $\boxtimes$ \\
JavaScript sayfada anlamlı bir işlev görüyor        & $\boxtimes$ \\
Sayfa telefonda da düzgün görünüyor                 & $\boxtimes$ \\
\bottomrule
\end{tabular}
\end{center}

% ====================================================================
\section{Masalların Seçimi}
% ====================================================================
Kurnaz Kurt: Kurnazlık temasını hayvan karakterler üzerinden işlemesi nedeniyle seçtim. Karakterlerin davranışları ve bu davranışların sonuçları, okuyucunun olayları değerlendirmesine olanak sağlıyor. Hayvanların yer aldığı bir anlatıya yer vererek derlemedeki karakter çeşitliliğini artırmayı amaçladım.

Sihirli Tavşan: Sihirli olaylar etrafında gelişen macerası ve hayal gücünü harekete geçiren yapısı nedeniyle seçtim. Olağanüstü bir karakterin hikâyeyi yönlendirmesi, sitenin masalsı atmosferini destekliyor. Tavşan karakteri ayrıca diğer masallardan ayırt edilebilen bir kart görseli oluşturmama imkân verdi.

Altın Saçlı Kız: Açgözlülüğün zararlarını ele alması ve karakterin seçimlerinin sonuçlarını göstermesi nedeniyle seçtim. Bu tema, okuyucunun istekler ve davranışlar arasındaki ilişki üzerine düşünmesini sağlıyor. Masalın belirgin karakteri, hikâyenin kart üzerinde kolayca tanınmasına da yardımcı oluyor.

Limon Kız: Olağanüstü olayları ve merak uyandıran adı nedeniyle seçtim. Hikâyenin hayal gücüne açık yapısı, ziyaretçiyi masalın devamını öğrenmeye yönlendiriyor. Limonla ilişkilendirilen karakter ve renkler, sitenin genel paleti içinde farklı bir görsel vurgu oluşturmama da olanak sağladı.

Akıllı Ormancı: Aklın güçten üstünlüğünü işlemesi ve sorunlara düşünerek çözüm bulmayı öne çıkarması nedeniyle seçtim. Karakterin karşılaştığı durumları nasıl değerlendirdiği, masalın temel mesajını anlaşılır kılıyor. Orman ortamı da sitenin doğadan ilham alan görsel diliyle ilişki kuruyor.

Altın Bülbül: Macera ve merak duygusunu canlı tutan anlatımı nedeniyle seçtim. Hikâyenin okuyucuyu bir sonraki olayı öğrenmeye yönlendirmesi, derlemeye hareketli bir okuma seçeneği katıyor. Altın bülbül figürü, masalın görsel olarak tanınmasını sağlayan bir odak noktası oluşturuyor.

Zümrüdüanka Kuşu: Umut ve iyilik temalarını olağanüstü bir kuş figürü üzerinden ele alması nedeniyle seçtim. Fantastik yapısı, sitenin sihirli atmosferini güçlendirirken okuyucunun hayalinde farklı bir dünya kurmasına yardımcı oluyor. Kuşun görkemli görünümü de illüstrasyonlarda dikkat çeken bir karakter oluşturmamı sağladı.

Gülen Ayva Ağlayan Nar: Cesaret ve empati temalarını işlemesi, ayrıca alışılmadık adıyla merak uyandırması nedeniyle seçtim. Gülme ve ağlama gibi insana ait duyguların meyvelerle ilişkilendirilmesi, masala özgün bir anlatım katıyor. Bu özellik, hikâyeyi hem içerik hem de görsel tasarım açısından derlemedeki diğer masallardan ayırıyor.

\newpage

\begin{longtable}{@{}cP{4cm}P{3cm}P{5.4cm}@{}}
\toprule
\textbf{No} & \textbf{Masal} & \textbf{Tür} & \textbf{Seçilme nedeni} \\
\midrule
\endhead
1 & Kurnaz Kurt & Hayvan Masalı & Kurnazlık teması \\
2 & Sihirli Tavşan & Olağanüstü masal & Sihirli bir macera. \\
3 & Altın Saçlı Kız & Olağanüstü masal & Açgözlülüğün zararını göstermesi. \\
4 & Limon Kız & Olağanüstü masal & Olağanüstü olaylarıyla merak uyandırdı.\\
5 & Akıllı Ormancı & Hayvan masalı &  Aklın güçten üstünlüğünü işlemesi.  \\
6 & Altın Bülbül & Olağanüstü masal & Macera ve merak duygusu uyandırdı. \\
7 & Zümrüdüanka Kuşu & Fantastik Anlatı & Umut ve İyiliği konu alması.\\
8 & Gülen Ayva Ağlayan Nar & Fantastik Anlatı & Cesaret ve empatiyi konu alması. \\
\bottomrule
\end{longtable}

% ====================================================================
\section{Tasarım Kararları}
% ====================================================================
Masal temasını yansıtmak için orman, gece ve ışık
unsurlarının yer aldığı bir görsel dil tercih ettim.
Her masalın içeriğine uygun illüstrasyonlar kullandım.
Küçük hareketli ışık noktaları, Masalsı Pırıltı adını
desteklerken ziyaretçiyi masalları keşfetmeye yönlendiriyor.
Kitap ve pırıltı simgelerini birleştiren logo da sitenin
okuma amacını yansıtıyor.

\subsection{Konuya Uygunluk}
Türk masallarındaki orman, saray ve olağanüstü varlıklardan
yola çıkarak masalsı bir atmosfer oluşturdum. Lacivert ve
mor tonlarını geceyi ve gizemi, turkuaz rengini ise ışıkları ve sihirli olayları çağrıştırması için seçtim. Başlıklarda kullanılan Georgia yazı tipi, basılı masal kitaplarını hatırlatıyor. İllüstrasyonlar ise her masalın karakterlerini ve geçtiği ortamı yansıtarak okuyucunun hikâyeyi hayal etmesini kolaylaştırıyor.

\subsection{Renk Paleti ve Yazı Tipi}
\begin{center}
\rowcolors{2}{acik}{white}
\begin{tabular}{@{}llP{6.4cm}@{}}
\toprule
\rowcolor{ana}\textcolor{white}{\textbf{Ad}} & \textcolor{white}{\textbf{Kod}} & \textcolor{white}{\textbf{Nerede / neden kullandınız?}} \\
\midrule
Lacivert & \texttt{\#0B1630} & Gece atmosferi için arka planda. \\
Mor & \texttt{\#B7A0E6} & Masalsı etki ve vurgu için. \\
Krem & \texttt{\#FFFAF2} & Okuma alanının zemininde. \\
Georgia / Arial & -- & Başlıklarda masalsı görünüm, arayüzde kolay okunabilirlik için. \\
\bottomrule
\end{tabular}
\end{center}

\subsection{Okunabilirlik, Gezinme ve Duyarlı Tasarım}
\begin{itemize}

\item \textbf{Okunabilirlik.}
Masal metinlerinin rahat okunabilmesi için okuma alanını
kart listesinden ayrı bir pencere olarak düzenledim.
Metnin satır uzunluğunu sınırlandırdım ve satır aralıklarını
geniş tutarak uzun paragrafların sıkışık görünmesini önledim.
Açık krem zemin üzerinde koyu metin kullanarak yazının
arka plandan kolayca ayrılmasını amaçladım.
Okuma ekranında yazı boyutu başlangıçta 20 piksel olarak
belirlendi. Kullanıcı, büyütme ve küçültme düğmeleriyle
bu değeri 16 ile 30 piksel arasında değiştirebiliyor.
Gece modu seçeneği de okuma alanının koyu bir görünüme
geçmesini sağlıyor. Böylece okuyucu, metnin görünümünü
kendi tercihine göre ayarlayabiliyor.

\item \textbf{Gezinme.}
Ziyaretçilerin masallara farklı yollarla ulaşabilmesi için
kartlar, arama kutusu ve tema filtrelerini birlikte kullandım.
Kartlarda masalın görseli, başlığı ve kısa tanıtımı yer alıyor.
Bu bilgiler, okuyucunun masalı açmadan önce seçim yapmasına
yardımcı oluyor. Arama kutusu başlığa göre sonuçları
daraltırken Sihir, Cesaret, İyilik ve Kurnazlık filtreleri
masalları temaları üzerinden keşfetmeyi sağlıyor.
Karar veremeyen ziyaretçiler için rastgele masal seçen
bir düğme de ekledim. Okuma ekranındaki önceki ve sonraki
masal düğmeleriyle kart listesine dönmeden farklı
hikâyelere geçilebiliyor. İlk ve son masalda ilgili geçiş
düğmesinin devre dışı kalması, listenin sınırlarını
kullanıcıya gösteriyor.

\item \textbf{Duyarlı tasarım.}
Sayfanın farklı ekran genişliklerine uyum sağlaması için
CSS medya sorguları kullandım. Geniş ekranlarda kartlar
yan yana yerleşirken ekran daraldığında sütun sayısı
azalıyor. Böylece kartların içindeki görsellerin,
başlıkların ve düğmelerin dar bir alana sıkışmasını
önlemeyi amaçladım. Menü, giriş bölümü ve okuma alanındaki
yerleşim ve boşluklar da ekran genişliğine göre değişiyor.
Mobil düzeni değerlendirirken yalnızca görünümü değil,
arama kutusuna ulaşmayı ve masalı açmayı da dikkate aldım.
Bu özelliklerin telefondaki davranışını ayrıca kontrol
etmek, masaüstü tasarımının mobilde nasıl kullanılacağını
daha iyi anlamamı sağladı.

\newpage

\item \textbf{Tutarlılık.}
Sayfanın farklı bölümlerinde ortak renkler, yazı tipleri,
kart biçimleri ve düğme stilleri kullanarak görsel
bütünlük oluşturdum. Başlıklarda Georgia, genel arayüzde
Arial yazı tipinden yararlandım. Kartların içerikleri
değişse de başlık, kısa açıklama ve okuma düğmesinin
benzer sırada bulunmasını sağladım.
Böylece ziyaretçi bir kartın düzenini öğrendiğinde
diğer kartlarda da bilgiyi kolayca bulabiliyor.
Renkleri ve boşlukları ortak CSS kurallarıyla yönetmek,
tasarımda yaptığım değişiklikleri ilgili alanlara
birlikte uygulamamı kolaylaştırdı.
Kartların üzerine gelindiğinde oluşan hareket ve gölge
değişiklikleri de benzer etkileşimlerde aynı görsel
geri bildirimin verilmesini sağlıyor.

\end{itemize}

\subsection{Ekran Görüntüleri}

\begin{center}
  \haftagorsel{ekran-mobil.png}

  \small Sayfanın telefon görünümü
\end{center}

\begin{center}
  \haftagorsel{ekran-masaustu.png}

  \small Sayfanın bilgisayar görünümü
\end{center}

\begin{center}
  \haftagorsel{okuma ekranı.png}

  \small Masal okuma ekranı
\end{center}



% \begin{figure}[h]
%   \centering
%   \includegraphics[width=0.9\textwidth]{ekran-masaustu.png}
%   \caption{Sayfanın bilgisayar görünümü}
% \end{figure}

\clearpage
\section{Gerçekleştirme}

\subsection{HTML Yapısı}
Sayfa; logo ve menünün bulunduğu üst alan, giriş bölümü,
masal kartları, okuma penceresi ve alt alandan oluşuyor.
Üst alan için header, gezinme bağlantıları için nav,
bölümler için section ve alt alan için footer etiketleri
kullanıldı. Başlıklar h1 ve h2, açıklamalar p etiketleriyle
oluşturuldu. Masallar dialog etiketiyle açılan bir
pencerede okunuyor. İşlem düğmeleri için button,
arama kutusu için input kullanıldı

\subsection{İçsel CSS}
Renkler CSS değişkenleriyle tanımlandı. Masal kartları
Grid, menü ve butonlar Flexbox ile düzenlendi.
Medya sorgularıyla düzenin ekran genişliğine uyum
sağlaması sağlandı. Işık noktaları CSS animasyonuyla
hareket ettirildi.

\begin{lstlisting}[caption={İçsel CSS örneği}]
body { background: #0b1630;
    color: #e9e8fa;}
\end{lstlisting}

\subsection{İçsel JavaScript}
Masallar JSON biçiminde dosyanın içinde tutuluyor.
render() fonksiyonu arama ve tema seçimine göre kartları
listeliyor. openStory() seçilen masalı okuma penceresinde
açıyor. Butonlara tıklama ve arama kutusuna yazma olayları
addEventListener ile işlevlere bağlanıyor.

\begin{lstlisting}[caption={İçsel JavaScript'ten önemli bir bölüm}]

search.addEventListener('input', render);
\end{lstlisting}



\newpage


\subsection{Sayfanın Özellikleri}

% ====================================================================
\section{Yapay Zekâ Kullanımı}
% ====================================================================
Kod desteği, arama, filtreleme ve okuma ayarları ile HTML, CSS ve JavaScript'in bazı bölümlerinde yapay zekâ kullanıldı. Masal illüstrasyonlarında da ChatGPT görsel üretiminden yararlanıldı.

\begin{longtable}{@{}P{2.4cm}P{3.4cm}P{4.1cm}P{3.4cm}@{}}
\toprule
\textbf{Araç} & \textbf{Ne için kullandınız?} & \textbf{İstem (prompt) özeti} & \textbf{Sayfada nerede?} \\
\midrule
\endhead
ChatGPT & Kod oluşturma ve açıklama & Sekiz masallı, arama ve okuma ayarları bulunan bir site hazırlama. & Sayfanın genel yapısı, arama, filtreleme ve okuma ekranı. \\
ChatGPT & Görsel üretimi & Mavi, mor ve turkuaz tonlarında masallara uygun illüstrasyonlar oluşturma ve düzenleme. & Ana sayfa arka planı ve masal illüstrasyonları. \\
\bottomrule
\end{longtable}

\subsection{Anladığım ve Kendim Yaptığım Kısımlar}
Masalları ve kaynaklarını seçtim. Sitenin adını, renk
paletini ve görsel atmosferini belirledim. Başlıklar,
kart biçimleri, karakterlerin görünümü ve ışık efektleri
için değişiklikler istedim. Kodun ve illüstrasyonların
oluşturulmasında yapay zekâ desteği aldım.

Yapay zekâ desteğiyle oluşturulan kod üzerinde başlıkları,
açıklamaları ve masal tanıtımlarını düzenledim. Renk,
yazı boyutu ve boşluk değerlerini değiştirerek tasarımı
tercihlerime göre geliştirdim. Değişikliklerin etkisini
önizlemede kontrol ettim.

% 

\newpage 
\section{Haftalık Katkılar}
% ====================================================================


\begin{longtable}{@{}>{\raggedright\arraybackslash}p{1.7cm}%
                    >{\raggedright\arraybackslash}p{6.4cm}%
                    >{\centering\arraybackslash}p{6cm}@{}}
\toprule
\rowcolor{ana}\textcolor{white}{\textbf{Hafta}} & \textcolor{white}{\textbf{Yapılan çalışma}} & \textcolor{white}{\textbf{Görsel}} \\
\midrule
\endfirsthead
\toprule
\rowcolor{ana}\textcolor{white}{\textbf{Hafta}} & \textcolor{white}{\textbf{Yapılan çalışma}} & \textcolor{white}{\textbf{Görsel}} \\
\midrule
\endhead
\endfoot

\textbf{1. hafta}\newline [tarih] &
[Bu hafta yaptığınız çalışma.]\newline \textbf{Commit:} [WorldWideWeb projesini anlatan bir sayfa hazırladım. İçeriği HTML’de başlık, paragraf ve bağlantılarla düzenledim. CSS ile yazı tipini, arka planı ve bağlantı renklerini belirledim.]\katilimkutulari &
\haftagorsel{hafta1.png}\haftaek{hafta1b.png} \\

\textbf{2. hafta}\newline [tarih] &
[Bu hafta yaptığınız çalışma.]\newline \textbf{Commit:} [Filmin renk paletini, simetrisini ve tipografisini anlatan bir web sayfası hazırladım. İçeriği ana ve alt başlıklarla ayırdım. Filmin görsel dünyasına uygun olarak mor arka plan ve pembe alt başlıklar kullandım.]\katilimkutulari &
\haftagorsel{hafta2.png}\haftaek{hafta2b.png} \\

\textbf{3. hafta}\newline [tarih] &
[Bu hafta yaptığınız çalışma.]\newline \textbf{Commit:} [Sayfaya “Günün Mesajı” adlı bir düğme ekledim. Düğmeye tıklandığında JavaScript ile boş paragrafın içeriğini güncelleyerek “Bugün güzel bir gün!” mesajını gösterdim.]\katilimkutulari &
\haftagorsel{hafta3.png}\haftaek{hafta3b.png} \\

\textbf{4. hafta}\newline [tarih] &
[Bu hafta yaptığınız çalışma.]\newline \textbf{Commit:} [İlk uygulamada, tıklanınca renkleri açılan dört kartlı bir eşleştirme oyunu hazırladım. JavaScript ile aynı renkteki kartların gizlenmesini, farklı renkteki kartların ise tekrar griye dönmesini sağladım.]\katilimkutulari &
\haftagorsel{hafta4.png}\haftaek{hafta4b.png} \\

\textbf{5. hafta}\newline [tarih] &
[Bu hafta yaptığınız çalışma.]\newline \textbf{Commit:} [CodePen’de düğmeye tıklandığında bir kutunun rengini değiştiren örnek üzerinde çalıştım. JavaScript’te if/else koşuluyla kutunun renginin gri ve sarı arasında nasıl değiştiğini inceledim.]\katilimkutulari &
\haftagorsel{hafta5.png}\haftaek{hafta5b.png} \\

\textbf{6. hafta}\newline [tarih] &
[Bu hafta yaptığınız çalışma.]\newline \textbf{Commit:} [CodePen’de “Renk Değiştir” düğmesine tıklanınca bir kutunun arka plan rengini değiştiren örnek üzerinde çalıştım. JavaScript’te if/else kullanarak kutu griyse sarıya, sarıysa griye dönmesini sağlayan koşulu inceledim.]\katilimkutulari &
\haftagorsel{hafta6.png}\haftaek{hafta6b.png} \\

\end{longtable}

% ====================================================================
\section{Test ve Değerlendirme}
% ====================================================================

\begin{center}
\rowcolors{2}{acik}{white}
\begin{tabular}{@{}P{6cm}P{3.4cm}P{3.4cm}@{}}
\toprule
\rowcolor{ana}\textcolor{white}{\textbf{Özellik}} & \textcolor{white}{\textbf{Bilgisayar}} & \textcolor{white}{\textbf{Telefon}} \\
\midrule
\relax Masal arama ve rastgele masal seçimi & Çalışıyor & Çalışıyor \\
\relax Masal okuma ekranının açılması & Çalışıyor & Çalışıyor \\
\relax Yazı boyutu değiştirme ve gece modu & Çalışıyor & Çalışıyor \\
\bottomrule
\end{tabular}
\end{center}


% ====================================================================
\section{Sunum Videosu}
% ====================================================================
\begin{center}
\begin{tabular}{@{}ll@{}}
\textbf{Video süresi:} & 12:13 \quad \aciklama{12--18 dakika olmalıdır} \\
\textbf{Görünürlük:} & [Listelenmemiş / Herkese açık] \\
\textbf{Bağlantı:} & \url{https://www.youtube.com/watch?v=1Jkv4_qK31A} \\
\end{tabular}
\end{center}

\subsection{Video Kontrol Listesi}
\aciklama{Teslimden önce hepsini kontrol edin ve kutuları $\boxtimes$ ile işaretleyin.}
\begin{itemize}
  \item[] $\boxtimes$ Süre 12--18 dakika aralığında.
  \item[] $\square$ Sunum boyunca kameram açık, yüzüm görünüyor.
  \item[] $\boxtimes$ Görünürlük \textbf{Listelenmemiş} ya da Herkese açık; \textbf{Özel} değil.
  \item[] $\boxtimes$ Bağlantıyı gizli pencerede açtım, video oynuyor.
  \item[] $\boxtimes$ YouTube videoyu işlemeyi bitirdi.
  \item[] $\square$ Videoyu baştan sona izledim; ses ve kamera tüm sunumda var.
  \item[] $\boxtimes$ Kapaktaki bağlantı ile bu bölümdeki bağlantı aynı.
\end{itemize}


\newpage
% ====================================================================
\section{Sonuç ve Kendi Değerlendirmem}
 

Masalsı Pırıltı projesinde sekiz masalı görseller ve etkileşimli özelliklerle bir araya getirerek bir web sayfasının geliştirilme sürecini deneyimledim. HTML ile içerik yapısını, CSS ile görünümü ve JavaScript ile kullanıcı işlemlerini nasıl yönetebileceğimi öğrendim. Renk, yazı boyutu ve yerleşim kararlarının sayfanın görünümü kadar kullanımını da etkilediğini fark ettim. En çok bazı özelliklerin telefonda ve bilgisayarda farklı çalışması konusunda zorlandım. Sorunun hangi aşamada ortaya çıktığını anlamak için arama ve masal açma işlemlerini ayrı ayrı denedim, düzenlemelerden sonra aynı adımları yeniden kontrol ettim. Yapay zekâ desteğiyle oluşturulan kodu incelemek ve değişikliklerin sonucunu önizlemede görmek, kod ile tasarım arasındaki ilişkiyi anlamama yardımcı oldu. Benzer bir projede mobil kontrolleri başlangıçtan itibaren yapar ve karşılaştığım sorunlarla çözümlerini daha düzenli kaydederim. İleride farklı kullanıcıların görüşlerini de alarak gezinmeyi ve okuma ekranını geliştirmek, teknik bilgimi artırarak kod üzerinde daha bağımsız çalışabilmek isterim.

% ====================================================================
\section{Kaynakça}



\begin{thebibliography}{9}
\raggedright

\bibitem{k1}
Türkiye Kültür Portalı.
\emph{Kurnaz Kurt Masalı}.
\url{https://www.kulturportali.gov.tr/portal/kurnazkurtmasali}

\bibitem{k2}
Masal.org.
\emph{Sihirli Tavşan}.
\url{https://www.masal.org/sihirli-tavsan/}

\bibitem{k3}
Masaloku.net.
\emph{Altın Saçlı Kız Masalı}.
\url{https://www.masaloku.net/altin-sacli-kiz-masali/}

\bibitem{k4}
Masal.org.
\emph{Limon Kız}.
\url{https://www.masal.org/limon-kiz/}

\bibitem{k5}

\emph{Akıllı Ormancı Masalı}.
\url{https://masal6.com/masal/akilli-ormanci-masalini-oku/}

\bibitem{k6}
Muzip Masal Cini.
\emph{Altın Bülbül Türk Masalı}.
\url{https://muzipmasalcini.com/altin-bulbul-turk-masali/}

\bibitem{k7}
Masal-oku.tr.
\emph{Zümrüdüanka Kuşu Masalı}.
\url{https://masal-oku.tr/zumruduanka-kusu-masali/}

\bibitem{k8}
Masaloku.net.
\emph{Gülen Ayva Ağlayan Nar}.
\url{https://www.masaloku.net/gulen-ayva-aglayan-nar/}

\end{thebibliography}

\newpage

% ====================================================================
\appendix
\section{Teslim Kontrol Listesi}
% ====================================================================
\aciklama{Teslim etmeden önce işaretleyin. Bu sayfayı teslim edilen PDF'te bırakabilirsiniz.}
\begin{itemize}
  \item[] $\boxtimes$ MEBİS'e \textbf{yalnızca bu raporun PDF'ini} yükleyeceğim. ZIP ve HTML dosyası yüklemeyeceğim.
  \item[] $\boxtimes$ PDF dosya adı \texttt{OgrenciNo\_AdSoyad\_vize.pdf} biçiminde.
  \item[] $\boxtimes$ Kapaktaki GitHub ve YouTube bağlantıları doğru; PDF'te tıklayıp denedim.
  \item[] $\boxtimes$ GitHub deposunda web sayfası (tek \texttt{.html}), rapor PDF'i, \texttt{.tex} kaynakları ve \texttt{README.md} var.
  \item[] $\boxtimes$ GitHub deposu \textbf{herkese açık} ya da öğretim üyesine erişim verdim.
  \item[] $\square$ Çalışmamın son hâli teslim saatinden \textbf{önce} depoya gönderildi.
  \item[] $\boxtimes$ Sayfa içsel CSS ve içsel JavaScript kullanıyor, harici dosya gerektirmiyor.
  \item[] $\boxtimes$ Hazır taslak kullanmadım; yapay zekâ kullandıysam Bölüm 5'i doldurdum.
  \item[] $\boxtimes$ Çalışmanın mantığını sunumda anlatabilirim.
\end{itemize}

\end{document}